% ============================================================
%  STANDALONE NOTE — Shortest-side skeleton theorem
%  ASTRA Lemma 6, rewritten in the notation of the current paper
%  (NOTATION LOCK v1.0). References "Lemma 3.4" etc. are to the
%  current integrated draft (CONTENT FROZEN 2026-09-27, now under
%  review). Nothing in this note is used by the paper yet.
% ============================================================
\documentclass[11pt]{article}
\usepackage{amsmath,amssymb,amsthm}
\usepackage[margin=1in]{geometry}
\usepackage{booktabs}
\theoremstyle{plain}
\newtheorem{theorem}{Theorem}
\newtheorem{lemma}[theorem]{Lemma}
\newtheorem{corollary}[theorem]{Corollary}
\theoremstyle{remark}
\newtheorem{remark}[theorem]{Remark}
\newcommand{\sm}{\setminus}
\DeclareMathOperator{\exc}{exc}
\title{Shortest-side skeleton theorem\\[4pt]
\large ASTRA Lemma 6 in the notation of the current paper --- audit copy}
\author{}
\date{2026-09-27 (rev 2: three editorial patches after independent audit)}
\begin{document}
\maketitle

\section*{0. Purpose and scope}

This note states and proves, in the notation of the current paper, the lemma that the
GPT-ASTRA run produced as its ``Lemma 6'': removing any shortest extreme line from an
edge-perimeter minimizer leaves an edge-perimeter minimizer whose threshold has dropped by
exactly one. It records precisely which statements of the current paper are used, which
step is new, and how the two exceptional residues coincide with the existing
Plateau-Start Bound. It is written to be audited on its own.

Every quantity is that of the current paper: \(p\), \(\sigma=p/2\), \(\lambda\), \(R\),
\(g(N)=\lceil\sqrt{12N-3}\rceil\), \(M(t)=\lfloor(t^{2}+3)/12\rfloor\),
\(N_{b}(\gamma)=\min\{N:g(N)=\gamma\}\), the tight hull \(H(S)=H(L;u,v,w)\), its side lines
\(\Lambda_{1},\dots,\Lambda_{6}\) and side lengths \(\ell_{i}\), the support traces
\(r_{i}=|S\cap\Lambda_{i}|\) of a directionally convex \(S\), and \(\rho(S)=\min_{i}r_{i}\).

\section*{1. Statements of the current paper that are used}

\begin{center}
\begin{tabular}{@{}llp{8.6cm}@{}}
\toprule
label & where & content used \\
\midrule
(P1) & Lemma 2.6 & \(p(S)=2R(S)\ge2\lambda(S)\); equality iff directionally convex \\
(P2) & Lemma 3.3 & \(S\subseteq H(S)\), \(\lambda(S)=\lambda(H)\), all six side lines meet \(S\) \\
(P3) & Lemma 3.4(ii),(iii) & \(H\) directionally convex with \(\sigma(H)=\lambda(H)\); \(\sum_{i}\ell_{i}=\lambda(H)+3\) \\
(P4) & Lemma 4.1(i),(ii) & peeling: \(S\sm\Lambda\) directionally convex polyhex with \(\sigma\) down by \(\ge1\); the removed cells can be restored one at a time through directionally convex polyhexes of \(\sigma\le\sigma(S)\) \\
(P5) & Thm 4.4(c),(d) & \(p_{\min}(N)=2g(N)\); a minimizer is directionally convex with \(\sigma=\lambda=g(N)\) \\
(P6) & definition of \(g\) & \(g(N)\le t\iff N\le M(t)\); \(g\) nondecreasing \\
(P7) & Lemma 4.6 & \(N_{b}(\gamma)=M(\gamma-1)+1\) \\
(P8) & Prop 11.1 & Plateau-Start Bound: for \(\gamma\ge7\) in the range of \(g\) and a minimizer \(S\) with \(|S|=N_{b}(\gamma)\), \(\rho(S)\le N_{b}(\gamma)-N_{b}(\gamma-1)\) \\
\bottomrule
\end{tabular}
\end{center}

(P8) is the only statement taken from Part~IV. Its own proof uses, besides (P2)--(P7): the
hull cell count of Lemma 3.4(iv) (to compute \(|H|\) for the regular hexagon and for the
hexagon with one opposite pair lengthened); the fact that the threshold jump
\(j(N)=g(N)-g(N-1)\) lies in \(\{0,1\}\) for \(N\ge3\) (to see that \(\gamma-1\) is in the
range of \(g\); in the present note this is replaced by the positivity of
\(M(\sigma-1)-M(\sigma-2)\) in Step 4); the Opposite-Pair Excess Lemma 11.2; and the
complete-boundary Lemma 11.3. It treats the residues \(\gamma\equiv3\) and
\(\gamma\equiv5\pmod 6\) as its Cases~2 and~3. In a one-proof reorganization these items
would move to the hull section of Part~I.

\section*{2. The theorem}

\begin{theorem}[shortest-side skeleton]\label{thm:skel}
Let \(S\) be an edge-perimeter minimizer with \(n=|S|\ge3\), let \(\sigma=\sigma(S)=g(n)\),
and let \(\Lambda\) be \emph{any} extreme line of \(S\) with \(|S\cap\Lambda|=\rho(S)\). Put
\(S':=S\sm\Lambda\). Then
\begin{enumerate}
\item[(i)] \(S'\) is nonempty and is an edge-perimeter minimizer with \(g(|S'|)=\sigma-1\);
\item[(ii)] the cells of \(S\cap\Lambda\) can be ordered \(c_{1},\dots,c_{\rho}\) so that each
\(T_{j}:=S'\cup\{c_{1},\dots,c_{j}\}\) is a directionally convex polyhex with
\(\exc(T_{j})\le2\).
\end{enumerate}
In words: peeling a shortest extreme line from a minimizer lands on a minimizer one
threshold step down, and putting the line back one cell at a time never costs more than
two units of edge perimeter.
\end{theorem}

\section*{3. Proof}

By (P5), \(S\) is directionally convex with \(\lambda(S)=\sigma=g(n)\), so
\(r_{1},\dots,r_{6}\) and \(\rho=\rho(S)\) are defined, and \(\rho\ge1\) by (P2). Let
\(H=H(S)\) with side lengths \(\ell_{i}\).

\medskip\noindent\textbf{Step 1 (upper bound on \(\rho\)).}
By (P2),(P3), \(\sum_{i}\ell_{i}=\lambda(H)+3=\sigma+3\). Since \(S\subseteq H\) we have
\(r_{i}\le\ell_{i}\), hence \(\sum_{i}r_{i}\le\sigma+3\). The minimum of six integers is at
most their mean, so
\begin{equation}\label{eq:rhobound}
   \rho\ \le\ \Bigl\lfloor\frac{\sigma+3}{6}\Bigr\rfloor .
\end{equation}
In words: the shortest support trace has at most one sixth of \(\sigma+3\) cells.

\medskip\noindent\textbf{Step 2 (lower bound on \(n\)).}
By (P6), \(g(n)=\sigma\) means \(n>M(\sigma-1)\), i.e.
\begin{equation}\label{eq:nbound}
   n\ \ge\ M(\sigma-1)+1 .
\end{equation}

\medskip\noindent\textbf{Step 3 (the residue table).}
Write \(\sigma=6m+s\) with \(0\le s\le5\). For \(t=6m+j\), \(0\le j\le5\), one has
\(M(t)=3m^{2}+mj+\lfloor(j^{2}+3)/12\rfloor\) with
\(\lfloor(j^{2}+3)/12\rfloor=0,0,0,1,1,2\) for \(j=0,\dots,5\). Taking differences (for
\(s\le1\) the arguments \(\sigma-1,\sigma-2\) fall into the block \(6(m-1)+j\)), and
computing \(\lfloor(\sigma+3)/6\rfloor=m+\lfloor(s+3)/6\rfloor\):

\begin{center}
\begin{tabular}{crr r}
\toprule
\(s=\sigma\bmod 6\) & \(M(\sigma-1)-M(\sigma-2)\) & \(\lfloor(\sigma+3)/6\rfloor\) & difference \\
\midrule
0 & \(m\)   & \(m\)   & \(0\) \\
1 & \(m\)   & \(m\)   & \(0\) \\
2 & \(m\)   & \(m\)   & \(0\) \\
3 & \(m\)   & \(m+1\) & \(\mathbf{-1}\) \\
4 & \(m+1\) & \(m+1\) & \(0\) \\
5 & \(m\)   & \(m+1\) & \(\mathbf{-1}\) \\
\bottomrule
\end{tabular}
\end{center}

Combining \eqref{eq:rhobound} and \eqref{eq:nbound},
\[
   n-\rho\ \ge\ M(\sigma-1)+1-\Bigl\lfloor\frac{\sigma+3}{6}\Bigr\rfloor .
\]
If \(s\notin\{3,5\}\) the table gives \(M(\sigma-1)-\lfloor(\sigma+3)/6\rfloor\ge M(\sigma-2)\),
hence
\begin{equation}\label{eq:strict}
   n-\rho\ \ge\ M(\sigma-2)+1 .
\end{equation}
If \(s\in\{3,5\}\) the table gives only \(n-\rho\ge M(\sigma-2)\), with equality if and only
if both \eqref{eq:nbound} and \eqref{eq:rhobound} are equalities:
\[
   n=M(\sigma-1)+1,\qquad \rho=m+1 .
\]
In words: the strict inequality \eqref{eq:strict} already holds except possibly when all
three of the following occur at once --- the threshold is \(3\) or \(5\) modulo \(6\), the
configuration is at the very start of its plateau, and the shortest trace is as long as
the counting allows.

\medskip\noindent\textbf{Step 4 (the two exceptional residues are the Plateau-Start Bound).}
Suppose \(s\in\{3,5\}\), \(n=M(\sigma-1)+1\), \(\rho=m+1\). By (P7),
\(n=N_{b}(\sigma)\), so \(S\) is a minimizer at the start of the plateau of \(\gamma:=\sigma\).
Since \(g(3)=\lceil\sqrt{33}\rceil=6\) and \(g\) is nondecreasing, \(n\ge3\) gives
\(\sigma\ge6\); hence \(m\ge1\), and as \(s\in\{3,5\}\), \(\gamma=\sigma\ge9\ge7\).
Moreover \(\gamma=g(n)\) lies in the range of \(g\), and in these residues the table gives
\(M(\sigma-1)-M(\sigma-2)=m\ge1\), so \(M(\sigma-2)<M(\sigma-1)\) and some \(N\) has
\(g(N)=\sigma-1\); thus \(\sigma-1\) lies in the range of \(g\) as well and
\(N_{b}(\sigma-1)\) is defined by (P7). Thus (P8) applies and gives
\[
   \rho(S)\ \le\ N_{b}(\sigma)-N_{b}(\sigma-1)
   \ =\ M(\sigma-1)-M(\sigma-2)\ =\ m
\]
by (P7) and the table. This contradicts \(\rho=m+1\). Hence the equality case does not occur,
and \eqref{eq:strict} holds for every residue.

\begin{remark}
(P8) is exactly ASTRA's ``Lemma 5'' in disguise. Its Cases~2 and~3 show that a minimizer at a
plateau start with \(\gamma\equiv3\) (resp.\ \(5\)) mod \(6\) and \(\rho=m+1\) would have all
six support traces of length \(\ge m+1\), which forces the hull to be the regular hexagon of
side \(m+1\) (resp.\ that hexagon with one opposite pair lengthened by one), and by
Lemma 11.3 forces \(S=H\); but \(|H|-N_{b}(\gamma)=m\ge1\). ASTRA arrives at the same two
hull shapes and the same count deficit \(m\).
\end{remark}

\medskip\noindent\textbf{Step 5 (conclusion (i)).}
From \eqref{eq:strict} and (P6), \(g(n-\rho)\ge\sigma-1\). Since \(\rho\ge1\) and
\(n-\rho\ge M(\sigma-2)+1\ge1\), the set \(S'=S\sm\Lambda\) is nonempty, so (P4)(i) applies:
\(S'\) is a directionally convex polyhex with \(\sigma(S')\le\sigma-1\). Then by (P1),(P5)
\[
   2(\sigma-1)\ \le\ 2g(|S'|)\ =\ p_{\min}(|S'|)\ \le\ p(S')\ =\ 2\sigma(S')\ \le\ 2(\sigma-1),
\]
so all are equalities: \(S'\) is a minimizer and \(g(|S'|)=\sigma-1\). Nothing in Steps 1--5
depends on which shortest line \(\Lambda\) was chosen, so (i) holds for every shortest
extreme line.

\medskip\noindent\textbf{Step 6 (conclusion (ii)).}
By (P4)(ii) there is an ordering \(c_{1},\dots,c_{\rho}\) of \(S\cap\Lambda\) such that every
\(T_{j}\) is a directionally convex polyhex with \(\sigma(T_{j})\le\sigma\). As
\(|T_{j}|\ge|S'|\) and \(g\) is nondecreasing, \(g(|T_{j}|)\ge g(|S'|)=\sigma-1\). Therefore,
using (P1),(P5),
\[
   \exc(T_{j})=p(T_{j})-2g(|T_{j}|)=2\bigl(\sigma(T_{j})-g(|T_{j}|)\bigr)
   \le2\bigl(\sigma-(\sigma-1)\bigr)=2 .
\]
In words: during the restoration the perimeter never exceeds the starting value while the
threshold is already one step lower, so the excess is at most two. \qed

\section*{4. Consequences}

\begin{corollary}[skeleton]\label{cor:skel}
Every minimizer \(S\) with \(|S|\ge1\) admits a chain
\(S=T_{0}\supset T_{1}\supset\cdots\supset T_{k}\) of minimizers, each obtained from the
previous by deleting a shortest extreme line, ending at \(|T_{k}|\le2\); and
\(g(|T_{i+1}|)=g(|T_{i}|)-1\) for each step.
\end{corollary}

\begin{proof}
Apply Theorem~\ref{thm:skel}(i) repeatedly while \(|T_{i}|\ge3\). Sizes strictly decrease.
\end{proof}

\begin{corollary}[one-cell growth with excess at most two]\label{cor:growth}
The main theorem of the current paper follows: interpolate each step of
Corollary~\ref{cor:skel} by Theorem~\ref{thm:skel}(ii), handle \(|T_{k}|\le2\) directly
(a single cell has \(\exc=0\); a domino has \(\exc=0\) and one of its cells is a single cell),
and read the whole chain backwards. Adjacency of each added cell follows from connectedness
of each stage, as in the current proof of Theorem 1.1.
\end{corollary}

\section*{5. What is new relative to the current paper}

\begin{itemize}
\item \textbf{New:} the inequality \(n-\rho\ge M(\sigma-2)+1\) for minimizers (Steps 1--4),
obtained from \eqref{eq:rhobound}, \eqref{eq:nbound}, the residue table and (P8). The current
Key Lemma 6.1 gives the weaker \(n-\rho\ge M(\sigma-2)\) on the larger near-extremal class and
does not use (P8).
\item \textbf{New:} Step 6 bounds the excess during restoration by comparing \(\sigma\) with
\(g(|S'|)=\sigma-1\) directly, without Theorem 4.4(e) and without the near-extremal class.
\item \textbf{Not new:} everything else is (P1)--(P8) verbatim.
\item \textbf{Consequence for the current text:} the sentences ``peeling a minimizer need
not produce a minimizer'' (\S5), Remark 7.3, and the description of \(X_{31}\) as ``a dead
end for the primary proof's induction'' (\S1.3) are false or misleading for shortest-line
peeling, which is what the primary proof performs.
\end{itemize}

\section*{6. Hypotheses, and where each is used}

\begin{center}
\begin{tabular}{@{}lp{9.5cm}@{}}
\toprule
\(S\) a minimizer & Step 2 (\(n>M(\sigma-1)\)); Step 4 via (P8); Step 5 (directional convexity from (P5)) \\
\(n\ge3\) & Step 4 (\(m\ge1\), \(\gamma\ge7\)); Step 5 (\(S'\neq\emptyset\)) \\
\(\Lambda\) shortest & Step 3 (\eqref{eq:rhobound} bounds \(|S\cap\Lambda|\)); Step 4 (\(\rho\) is what (P8) bounds) \\
\bottomrule
\end{tabular}
\end{center}

The word ``shortest'' cannot be dropped: for a non-shortest line \(|S\cap\Lambda|\) may exceed
\(\lfloor(\sigma+3)/6\rfloor\), and the computation reported below shows that roughly half of
the non-shortest peelings of minimizers with \(N\le34\) do not land on a minimizer.

\section*{7. Machine check (re-check only, not part of the proof)}

Independent implementation (cube coordinates, no shared code with the paper's ancillary
files), see \texttt{skeleton\_check.py}:
\begin{itemize}
\item all minimizers with \(3\le N\le34\), all shortest extreme lines (ties included):
\(2{,}382\) peelings, every one a minimizer with \(g\) down by exactly \(1\), and every
restoration order of Lemma 4.1(ii) keeps \(\exc\le2\);
\item non-shortest peelings: \(6{,}498\), of which \(3{,}192\) land on a minimizer;
\item residue table verified for \(\sigma\le2999\); \(N_{b}(\sigma)=M(\sigma-1)+1\) and the
two boundary counts \(3m^{2}+2m+1\), \(3m^{2}+4m+2\) verified for \(m\le59\).
\end{itemize}

\end{document}
